\section{Structural limits}\label{sec:limits}

The algorithms of Sections~\ref{sec:grids}--\ref{sec:exact} rest on four
hypotheses: a product domain, upper bounds on coordinate curvature, weighted
or point growth, and a correction that charges for rounding every coordinate.
Section~\ref{sec:recourse} replaces the corrections of eliminated coordinates
by exact or certified value functions, Section~\ref{sec:constraints} replaces
the product domain by a feasible rounding law for totally unimodular
constraints, and Sections~\ref{sec:constraints} and~\ref{sec:optsets} use set
growth instead of growth at one minimizer. This section shows what fails when a
hypothesis is dropped without such a replacement, and that the dependence of
Theorem~\ref{thm:approx} on $p$ and $\kappa$ is close to necessary.
\begin{itemize}
\item \emph{Conditioning and width}
(Sections~\ref{sec:limits-conditioning}--\ref{sec:limits-width}). Unless
P${}={}$NP, neither $p$ nor $\kappa$ can be dropped, and the dependence on
$\kappa$ cannot be polylogarithmic (Corollary~\ref{lim:cor:nopolylog} and
Proposition~\ref{prop:lbwidth}). Under the exponential-time hypothesis (ETH)
the dependence on $p$ is exponential even for $\kappa\le2$
(Proposition~\ref{prop:lbwidth}); under randomized ETH the exponent of
$\kappa$ cannot be $o(p)$, already for integer box quadratics
(Proposition~\ref{prop:lbproduct}). In the value-oracle model the exponent
$p/2$ of $\kappa$ is optimal (Proposition~\ref{lim:prop:oracle}).
\item \emph{Exact messages} (Section~\ref{sec:limits-messages}). Exact
piecewise-quadratic messages can need exponentially many pieces at bag size
three and $\kappa\le80$ (Proposition~\ref{lim:prop:messages}).
\item \emph{Set growth} (Section~\ref{sec:limits-setgrowth}). Growth towards a
continuum of minimizers does not bound the grids produced by filtering
(Proposition~\ref{lim:prop:setgrowth}). A deterministic algorithm that
queries values and derivatives and is correct under set growth needs at least
$(1/(2\sqrt{6\varepsilon})-1)^n-1$ queries on the zero function in $n$
variables, for which $\kappa_S=1$ (Proposition~\ref{prop:oraclebarrier}).
With two isolated minimizers, the single-center graded grids of
Section~\ref{sec:growth} need exponential time
(Proposition~\ref{prop:twocenters} in Section~\ref{sec:optsets}).
\item \emph{Product domains} (Section~\ref{sec:limits-constraints}). A path of
local affine equalities makes the problem NP-hard at bag size three and
$\kappa=1$ (Proposition~\ref{lim:prop:constraints}).
\item \emph{Local corrections} (Section~\ref{sec:limits-local}). Corrections
charged to the coordinates of one bag are not valid lower bounds
(Proposition~\ref{prop:star} in Section~\ref{sec:recourse}), and matching
separator moments of any finite order does not give an error bound in terms
of negative curvature (Proposition~\ref{lim:prop:moments} in
Appendix~\ref{app:moments}).
\end{itemize}
Del Pia and Khajavirad prove that continuous box QP is strongly NP-hard at
treewidth two, even with bounded integer coefficients, and that quartic box
problems on paths are strongly NP-hard \cite{DelPiaKhajavirad2026}. Their
instances need not have a unique minimizer, so their result does not by
itself show that $\kappa$ must enter the running time.

\subsection{Conditioning cannot be dropped}\label{sec:limits-conditioning}

An instance of Subset Sum consists of positive integers $a_1,\ldots,a_m$ and
a target $a_0$; it asks whether some subset of $\{a_1,\ldots,a_m\}$ sums to
$a_0$. The following construction refines the simple weak NP-hardness
reduction in \cite[Remark~2]{DelPiaKhajavirad2026}: a stronger concave penalty
puts every minimizer at a vertex, a lexicographic linear term makes the
minimizer unique, and the growth constant is computed explicitly.

\begin{proposition}[Unique minimizers without conditioning]\label{lim:prop:unique}
Let $a_1,\ldots,a_m$ be positive integers, $m\ge2$, let $a=(a_1,\ldots,a_m)$,
and let $a_0$ be an integer with $0\le a_0\le A:=\sum_ia_i$. Put
$\lambda=1/(m2^{m+1})$. On the box
\[
x\in[0,1]^m,\qquad y=(y_1,\ldots,y_{m-1})\in[-A,2A]^{m-1},
\qquad y_0:=0,\quad y_m:=a_0,
\]
minimize
\[
F(x,y)=\sum_{i=1}^m(y_i-y_{i-1}-a_ix_i)^2
+\norm a^2\sum_{i=1}^mx_i(1-x_i)+\lambda\sum_{i=1}^m2^{i-1}x_i .
\]
Then:
\begin{enumerate}
\item The bags $\{y_{i-1},y_i,x_i\}$, with the constants $y_0,y_m$ omitted
  and arranged as a path in the order $i=1,\ldots,m$, form a tree
  decomposition with $p\le3$. All data have binary length polynomial in
  that of $(a,a_0)$.
\item $F$ has a unique minimizer $z^*=(x^*,y^*)$, and $x^*\in\{0,1\}^m$.
  If some $v\in\{0,1\}^m$ satisfies $a^{\T}v=a_0$, then
  $a^{\T}x^*=a_0$ and $\OPT<1/(2m)$; otherwise $\OPT\ge1/m$.
\item $L=4$ is a valid upper coordinate-curvature bound and
  \[
  g=\frac{\lambda}{m\,(1+2(m-1)\norm a^2)}
  \]
  is a valid growth constant. Hence
  $\kappa\le2^{m+3}m^2(1+2m\norm a^2)$ and $\log_2\kappa=O(I)$.
\end{enumerate}
\end{proposition}

\begin{proof}
Part 1 is immediate: $y_i$ lies in the consecutive bags $i$ and $i+1$, each
$x_i$ in one bag, and each square and unary term lies in a bag.

\emph{Elimination of $y$.} For $x\in[0,1]^m$ let
$e(x)=a^{\T}x-a_0$, $\sigma_0(x)=0$, and
\[
\sigma_k(x)=\sum_{j\le k}a_jx_j-\frac kme(x)\qquad(k=1,\ldots,m),
\]
so that $\sigma_m(x)=a_0$. For a feasible $y$ put $\delta_k=y_k-\sigma_k(x)$,
with $\delta_0=\delta_m=0$. Then
$y_i-y_{i-1}-a_ix_i=-e(x)/m+\delta_i-\delta_{i-1}$, and
$\sum_i(\delta_i-\delta_{i-1})=0$. Expanding the squares gives the identity
\begin{equation}\label{lim:eq:split}
F(x,y)=\Phi(x)+\sum_{i=1}^m(\delta_i-\delta_{i-1})^2,\qquad
\Phi(x)=\frac{e(x)^2}m+\norm a^2\sum_ix_i(1-x_i)+\lambda\sum_i2^{i-1}x_i .
\end{equation}
Because $0\le\sum_{j\le k}a_jx_j\le A$ and $\abs{e(x)}\le A$, the vector
$(\sigma_1(x),\ldots,\sigma_{m-1}(x))$ lies in $[-A,2A]^{m-1}$.
Since $\delta_k=\sum_{i\le k}(\delta_i-\delta_{i-1})$, the Cauchy--Schwarz
inequality gives $\delta_k^2\le k\sum_i(\delta_i-\delta_{i-1})^2$. Summing
over $k\le m-1$ gives the path inequality
\begin{equation}\label{lim:eq:path}
\sum_{i=1}^m(\delta_i-\delta_{i-1})^2\ge\frac2{m(m-1)}\sum_{k=1}^{m-1}
\delta_k^2\qquad(\delta_0=\delta_m=0),
\end{equation}
which holds for all reals $\delta_0,\ldots,\delta_m$ with $\delta_0=\delta_m=0$.

\emph{The reduced function.} Its Hessian is
$\nabla^2\Phi=\tfrac2maa^{\T}-2\norm a^2I\preceq0$, so $\Phi$ is concave. At a
vertex $v$, $\Phi(v)=e(v)^2/m+\lambda\chi(v)$ with
$\chi(v)=\sum_i2^{i-1}v_i$, a bijection of $\{0,1\}^m$ onto
$\{0,\ldots,2^m-1\}$. If $\Phi(v)=\Phi(v')$, then
$(e(v)^2-e(v')^2)/m=\lambda(\chi(v')-\chi(v))$; the left side lies in
$m^{-1}\Z$ and the right side has absolute value less than
$\lambda2^m=1/(2m)$, so both vanish and $v=v'$. Let $x^*$ be the unique
minimizing vertex and $\Phi^*=\Phi(x^*)$. We claim that
$\Phi(v)-\Phi^*\ge\lambda$ for every vertex $v\ne x^*$. If
$e(v)^2=e(x^*)^2$, the difference is a positive integer multiple of
$\lambda$. Otherwise $e(v)^2-e(x^*)^2$ is a nonzero integer. It cannot be
negative, since then $\Phi(v)-\Phi^*\le-1/m+\lambda(2^m-1)<0$. Hence it
is at least one and $\Phi(v)-\Phi^*\ge1/m-\lambda(2^m-1)>1/(2m)\ge\lambda$.

For $x\in[0,1]^m$ let $Y$ have independent coordinates with
$\Pr(Y_i=1)=x_i$. By concavity and Jensen's inequality,
$\Phi(x)=\Phi(\E Y)\ge\E\Phi(Y)\ge\Phi^*+\lambda\Pr(Y\ne x^*)$. Since $x^*$ is
binary, $\Pr(Y_i\ne x_i^*)=\abs{x_i-x_i^*}\le1$, so
$\Pr(Y\ne x^*)\ge\max_i\abs{x_i-x_i^*}\ge\frac1m\norm{x-x^*}^2$. Therefore
\begin{equation}\label{lim:eq:phigrowth}
\Phi(x)-\Phi^*\ge\frac\lambda m\norm{x-x^*}^2\qquad(x\in[0,1]^m).
\end{equation}

\emph{Uniqueness and growth.} By \eqref{lim:eq:split} and
\eqref{lim:eq:phigrowth}, $z^*=(x^*,\sigma(x^*))$ is the unique minimizer
and $\OPT=\Phi^*$. Let $z=(x,y)$ be feasible and $h=x-x^*$. Then
$\sigma_k(x)-\sigma_k(x^*)=\sum_jc_{kj}a_jh_j$ with
$c_{kj}=[j\le k]-k/m\in[-1,1]$, so
$\norm{\sigma(x)-\sigma(x^*)}^2\le(m-1)\norm a^2\norm h^2$ and
\[
\norm{z-z^*}^2\le\norm h^2+2\norm\delta^2+2\norm{\sigma(x)-\sigma(x^*)}^2
\le(1+2(m-1)\norm a^2)\norm h^2+2\norm\delta^2 .
\]
By \eqref{lim:eq:split}--\eqref{lim:eq:phigrowth},
$F(z)-\OPT\ge(\lambda/m)\norm h^2+\frac2{m(m-1)}\norm\delta^2$. Since
$g(1+2(m-1)\norm a^2)=\lambda/m$ and $g\le\lambda/m<1/(m(m-1))$, this gives
$F(z)-\OPT\ge g\norm{z-z^*}^2$. On coordinate lines,
$\partial^2F/\partial y_k^2=4$ and
$\partial^2F/\partial x_i^2=2a_i^2-2\norm a^2\le0$, so $L=4$ is valid and
$\kappa\le4/g\le2^{m+3}m^2(1+2m\norm a^2)$. Both $m$ and $\log\norm a^2$ are
$O(I)$.

\emph{Decision.} If $a^{\T}v=a_0$ for a vertex $v$, then
$\OPT\le\Phi(v)\le\lambda(2^m-1)<1/(2m)$, and, as shown above,
$e(x^*)^2\le e(v)^2=0$. If no vertex solves the instance, every vertex has
$e(v)^2\ge1$, and $\OPT=\Phi(x^*)\ge1/m$.
\end{proof}

\begin{corollary}[The dependence on $\kappa$ is not polylogarithmic]
\label{lim:cor:nopolylog}
Unless $\mathrm P=\mathrm{NP}$, no algorithm computes, for every rational
continuous box QP with a supplied tree decomposition of maximum bag size at
most three and a unique minimizer, the optimal value within $2^{-q}$ in time
polynomial in $I+q+\log\kappa$. In particular, no algorithm achieves the
guarantee of Theorem~\ref{thm:approx} with $f(3,\kappa)$ bounded by a
polynomial in $\log\kappa$.
\end{corollary}

\begin{proof}
Apply such an algorithm to the instances of
Proposition~\ref{lim:prop:unique} with $q=\lceil\log_2(4m)\rceil$. An
approximation within $1/(4m)$ separates $\OPT<1/(2m)$ from $\OPT\ge1/m$,
and $\log\kappa=O(I)$. This would decide Subset Sum, which is NP-complete
\cite{GareyJohnson1979}, in polynomial time.
\end{proof}

For fixed $p$, the function
$f(p,\kappa)=c_0(c_1p\sqrt\kappa)^p\kappa(1+\log_2\kappa)^2$ of
Theorem~\ref{thm:approx} is bounded by a polynomial in $\kappa$, and the proof
of Theorem~\ref{thm:exact} bounds $f_1(p,\kappa)$ by $f(p,\kappa)$ times a
power of $1+\log_2\kappa$. Hence, for fixed $p$, instances with a unique
minimizer and $\kappa\le I^{O(1)}$ are solved exactly in polynomial time,
while by Proposition~\ref{lim:prop:unique} the problem is NP-hard on
instances with $\kappa\le2^{O(I)}$. Adding decoupled terms $t_j^2$ with
$t_j\in[0,1]$, each in a bag of its own, keeps the minimizer unique, leaves
$p$ and $\kappa$ unchanged, and increases $I$ by an amount of our choice; so
for every fixed $\delta>0$ the problem is NP-hard already on instances with
$\kappa\le2^{I^\delta}$.

The strongly NP-hard instances of \cite[Theorem~3]{DelPiaKhajavirad2026}
have bags of size three and coefficients bounded by absolute constants.
Remark~\ref{lim:rem:dk} in Appendix~\ref{app:dk} shows that, already for a
family of two-item instances on which their minimizer is unique, $\kappa$ and the ratio $\nu/g$ of
negative curvature to growth are exponential in the input length. These
instances are therefore consistent with Corollary~\ref{lim:cor:nopolylog},
and they do not contradict an algorithm parameterized by $\kappa$ or by
$\nu/g$.

\subsection{The value-oracle model}\label{sec:limits-oracle}

CT (Algorithm~\ref{alg:ct}) uses only exact values of $F$ and the curvature
bounds. In a model where $F$ is available only through evaluations, the
exponent $p/2$ of $\kappa$ in its running time cannot be improved. The
construction is the classical hidden-well argument
\cite{NemirovskiYudin1983}, adapted to the present hypotheses.

\begin{proposition}[Value-oracle lower bound]\label{lim:prop:oracle}
Let $p\ge1$, $L>0$, $\kappa\ge8p$, $g=L/\kappa$, and $X=[0,1]^p$. For
$c\in X$ put
\[
F_c(x)=\min\Bigl\{gp,\ \frac L2\norm{x-c}^2\Bigr\}.
\]
Each $F_c$ has upper coordinate curvature $L$, the unique minimizer $c$
with $F_c(c)=0$, and growth $F_c(x)\ge g\norm{x-c}^2$ on $X$; a single
bag of size $p$ is a valid decomposition. Suppose a deterministic algorithm
receives $p$, $L$, $\kappa$ and value access to $F_c$ for an unknown $c$,
and always returns $y\in X$ with $F_c(y)<Lp/\kappa$. Then on some $F_c$ it
makes at least $(\kappa/(8p))^{p/2}-1$ evaluations. A randomized algorithm
that makes at most $Q$ evaluations succeeds, for some $c$, with probability
at most $(Q+1)(8p/\kappa)^{p/2}$.
\end{proposition}

\begin{proof}
On a coordinate line, $\frac L2\norm{x-c}^2-\frac L2x_i^2$ is affine in
$x_i$ and $gp-\frac L2x_i^2$ is concave; their minimum is concave. Thus
$t\mapsto F_c-Lt^2/2$ is concave on every coordinate line. For $x\in X$,
$\norm{x-c}^2\le p$, so $gp\ge g\norm{x-c}^2$; and $L/2\ge g$ because
$\kappa\ge2$. This proves growth, and $c$ is the unique zero.

Moreover $F_c(x)=gp$ unless $\norm{x-c}<r$, where $r^2=2gp/L=2p/\kappa$.
Since $\kappa\ge8p$, $2r\le1$. Let $\mathcal C$ consist of the points of
$X$ whose coordinates lie in $\{0,2r,4r,\ldots\}$. Then
$\abs{\mathcal C}=(\lfloor1/(2r)\rfloor+1)^p\ge(2r)^{-p}=(\kappa/(8p))^{p/2}$,
and the open balls $B(c,r)$, $c\in\mathcal C$, are pairwise disjoint. The
guarantee $F_c(y)<Lp/\kappa=gp$ means $y\in B(c,r)$.

Consider the run of the deterministic algorithm in which every answer equals
$gp$, with evaluation points $x^1,x^2,\ldots$ (finitely or infinitely many).
For $c\in\mathcal C$ let $t_c$ be the index of the first evaluation point in
$B(c,r)$, and $t_c=\infty$ if there is none. Since $F_c=gp$ outside $B(c,r)$,
the run on $F_c$ agrees with the constant run up to and including evaluation
$t_c$. If $t_c=\infty$, the two runs coincide, so the constant run stops and
its output lies in $B(c,r)$; as the balls are disjoint, this happens for at
most one $c$. Each point lies in at most one ball, so the finite $t_c$ are
distinct. At least $\abs{\mathcal C}-1$ of them are finite, so the largest
is at least $\abs{\mathcal C}-1$, and the run on the corresponding $F_c$
makes at least $\abs{\mathcal C}-1$ evaluations.

For a randomized algorithm, choose $c$ uniformly in $\mathcal C$ and fix the
random bits. Consider the first $Q$ evaluation points of the run on the
constant answers. If $B(c,r)$ contains none of them, the run on $F_c$ agrees
with the constant run on these evaluations; since it makes at most $Q$
evaluations, it stops there with the output of the constant run, and it
succeeds only if that output lies in $B(c,r)$. Hence at most $Q+1$ choices of
$c$ succeed. Averaging over the random bits gives success probability at most
$(Q+1)/\abs{\mathcal C}$ for a uniform $c$, hence for some $c$.
\end{proof}

\begin{remark}[Comparison with the algorithm]\label{lim:rem:oracle}
With one bag ($n=p$, $N=1$), each stage of CT evaluates $F$ at most at
$K^p$ grid points. By Lemma~\ref{lem:states} and the trial schedule,
$K=O(\sqrt\kappa\log(p+2))$ nodes per coordinate in the successful trial,
and the earlier trials cost no more. Thus the exponent $p/2$ of $\kappa$ is
attained, and Proposition~\ref{lim:prop:oracle} shows that it cannot be
improved in this model; the remaining gap is a factor
$(C\sqrt p\log(p+2))^p$ times the number of stages. The lower bound is also
exponential in $p$ once $\kappa\ge32p$.
\end{remark}

\subsection{Width and the exponent of \texorpdfstring{$\kappa$}{kappa}}\label{sec:limits-width}

For explicit quadratics we use the exponential-time hypothesis (ETH): for some
$c>0$, no algorithm decides 3-SAT on $n'$ variables in time $2^{cn'}$
\cite{ImpagliazzoPaturiZane2001}. The randomized exponential-time hypothesis
(rETH) states the same for randomized algorithms with error probability at
most $1/3$ \cite{DellEtAl2014}. The next construction is the standard box
formulation of maximum independent set, with a binary tie-break that makes
the minimizer unique.

\begin{proposition}[Exponential dependence on the width at bounded conditioning]
\label{prop:lbwidth}
For a graph with vertex set $[n]$, $n\ge1$, and edge set $E$ let
\[
\Psi(x)=2^n\Bigl(-\sum_ix_i+2\sum_{ij\in E}x_ix_j\Bigr)+\sum_i2^{i-1}x_i
+\frac1{2n}\sum_i(x_i^2-x_i),\qquad x\in[0,1]^n .
\]
Then $\Psi$ has a unique minimizer $x^*$, which is binary,
$\lfloor\Psi^*/2^n\rfloor=-\alpha$, where $\Psi^*=\min\Psi$ and $\alpha$ is
the independence number of the graph, and
$\Psi(x)-\Psi^*\ge\frac1{2n}\norm{x-x^*}^2$ on $[0,1]^n$. With $L_i=L=1/n$
this gives $\kappa\le2$ and $\bar\kappa\le2$. The same holds on
$\{0,1\}^n$. Every tree decomposition of the graph is a decomposition of
$\Psi$. Consequently, computing $\OPT$ within $1/2$ for box quadratics with
$\kappa\le2$ (continuous, or all-binary) is NP-hard, and, assuming ETH, no
algorithm does it in time $2^{o(p)}I^{O(1)}$.
\end{proposition}

\begin{proof}
On binary points $\Psi=2^n\Phi+\chi$ with $\Phi=-\sum x_i+2\sum_Ex_ix_j$
integral and $\chi=\sum2^{i-1}x_i\in\{0,\dots,2^n-1\}$ injective. A binary
minimizer therefore minimizes $\Phi$ and is unique. If the support of a
binary $x$ contains both endpoints of an edge, deleting such an endpoint $u$
from the support changes $\Phi$ by $1-2d(u)\le-1$, where $d(u)\ge1$ is the
number of neighbors of $u$ in the support; hence $\min\Phi=-\alpha$.

For $x\in[0,1]^n$ write $\Psi=\Psi_0+\frac1{2n}\sum(x_i^2-x_i)$, where
$\Psi_0=2^n\Phi+\chi$ is multilinear. Let $Y$ have independent coordinates
with $\Pr[Y_i=1]=x_i$. Then $\Psi_0(x)=\E\Psi_0(Y)=\E\Psi(Y)$, and since
$\Psi(v)\ge\Psi(x^*)+1$ for binary $v\ne x^*$,
$\Psi_0(x)-\Psi(x^*)\ge\Pr[Y\ne x^*]$. With
$\delta_i=\abs{x_i-x_i^*}\in[0,1]$,
$\Pr[Y=x^*]=\prod_i(1-\delta_i)\le1-\max_i\delta_i$, so
\begin{equation}\label{lim:eq:lbwidth}
\Psi_0(x)-\Psi(x^*)\ge\max_i\delta_i\ge\frac1n\sum_i\delta_i
\ge\frac1n\norm{x-x^*}^2 .
\end{equation}
Also $x_i(1-x_i)\le\delta_i$ because $x_i^*\in\{0,1\}$. Hence
$\Psi(x)-\Psi(x^*)\ge(\frac1n-\frac1{2n})\sum_i\delta_i\ge\frac1{2n}\norm{x-x^*}^2$,
so $x^*$ is the unique minimizer over $[0,1]^n$ and $\Psi^*=\Psi(x^*)$. The
only squared terms are $\frac1{2n}x_i^2$, so $L=1/n$ and $L/g\le2$; in the
weighted norm, $\Psi(x)-\Psi^*\ge\frac12\norm{x-x^*}_{\mathbf L}^2$, so
$\bar\kappa\le2$.

Since $n\ge1$, we have $\alpha\ge1$, so $\Phi(x^*)=-\alpha<\Phi(0)$ and
$x^*\ne0$. Thus $\Psi^*=-2^n\alpha+\chi(x^*)$ with
$\chi(x^*)\in[1,2^n-1]$, and every
number $\tilde\Psi$ with $\abs{\tilde\Psi-\Psi^*}\le1/2$ satisfies
$\lfloor\tilde\Psi/2^n\rfloor=-\alpha$, which
also gives $\lfloor\Psi^*/2^n\rfloor=-\alpha$. The instance has $O(n^2)$
coefficients of $O(n)$ bits, so $I=n^{O(1)}$, and the single bag $[n]$ gives
$p=n$. Computing $\alpha$ is NP-hard \cite{GareyJohnson1979}, and an
algorithm with running time $2^{o(p)}I^{O(1)}$ would compute it in time
$2^{o(n)}$, contradicting ETH \cite{ImpagliazzoPaturiZane2001}.
\end{proof}

Under the strong exponential-time hypothesis, the same reduction, applied with
a supplied tree decomposition of the graph, excludes running times
$(2-\epsilon)^pI^{O(1)}$ for every $\epsilon>0$, by the lower bound of
\cite{LokshtanovMarxSaurabh2018} for independent set.

If $L=0$, CT is exact at its first stage, whose grids consist of box
endpoints, and it runs in time $2^p\poly(I)$. The multilinear part $\Psi_0$
of the instance above has $L=0$, the same unique binary minimizer, and growth
constant $g=1/n$ by \eqref{lim:eq:lbwidth}, so $\kappa=1$; the same reduction
applies to it. For $L=0$ the dependence on $p$ is therefore optimal up to the
base of the exponential under ETH, and optimal under the strong
exponential-time hypothesis.

The next proposition bounds the exponent of $\kappa$. Its reduction uses
random weights, so it needs rETH.

\begin{proposition}[The exponent of $\kappa$ cannot be $o(p)$]
\label{prop:lbproduct}
Assume rETH. There are no computable $f$ and computable nondecreasing
unbounded $\psi\ge1$ such that some algorithm, on every integer box quadratic
with a supplied tree decomposition of maximum bag size $p$ and a unique
minimizer, outputs a feasible point within $1/2$ of $\OPT$ in time
$f(p)\,(\kappa I)^{p/\psi(p)}$. In particular, no algorithm has running time
$f(p)\kappa^{a(p)}I^C$ with a computable $f$, a constant $C$, and
$a(p)\le p/\psi(p)$ for such a $\psi$. The bound of
Theorem~\ref{thm:approx} has $a(p)=p/2+O(1)$.
\end{proposition}

The proof, a reduction from multicolored clique with randomized isolation,
is in Appendix~\ref{app:lbproduct}. Under rETH, the exponent $p/2+O(1)$ of
$\kappa$ in Theorem~\ref{thm:approx} therefore cannot be improved to $o(p)$,
in the sense $a(p)\le p/\psi(p)$ above, on integer instances. The constant in the
exponent, the necessity of the factor $p^p$, and whether the same holds for
purely continuous instances or with a deterministic reduction, are open.

\subsection{Exact messages can be exponentially large}
\label{sec:limits-messages}

For forests, exact value functions of box QP are piecewise quadratic with
linearly many pieces \cite{DelPiaKhajavirad2026}. Dynamic programming with
exact piecewise-quadratic messages on trees is also used for total-variation
problems \cite{KolmogorovPockRolinek2016,KuricAhmetspahicPock2024}; for
nonconvex piecewise-quadratic costs the worst-case bound of
\cite{KuricAhmetspahicPock2024} is exponential. The next proposition gives a
lower bound for the message itself: at bag size three, exact messages can have
exponentially many pieces even when $\kappa$ is bounded and the minimizer is
unique. This is why CT uses grids rather than exact piecewise-quadratic
messages.

\begin{proposition}[Exponential messages at bounded conditioning]
\label{lim:prop:messages}
For $m\ge2$ let $\xi_t\in[0,2^t-1]$ and $z_t\in[0,1]$, $t=1,\ldots,m$, put
$\xi_0:=0$, $\rho_t=\xi_t-2\xi_{t-1}-z_t$, and
\[
\Psi_m(\xi,z)=\xi_m^2+\sum_{t=1}^m\rho_t^2+\frac18\sum_{t=1}^mz_t(1-z_t).
\]
Then:
\begin{enumerate}
\item the bags $\{\xi_{t-1},\xi_t,z_t\}$ ($\xi_0$ omitted) form a path
  decomposition with $p=3$, and $I=O(m^2)$;
\item $\Psi_m(\xi,z)\ge\frac18(\norm\xi^2+\norm z^2)$ on the box, so zero is
  the unique minimizer and $g=1/8$ is valid;
\item $L=10$ is the exact maximum coordinate curvature, so $\kappa\le80$;
  moreover $\nabla^2\Psi_m\succeq-\frac14I$;
\item let $V(v)=\min\{\Psi_m(\xi,z):\xi_m=v\}$ for $v\in[0,2^m-1]$, the
  value function of $\xi_m$. Up to the term $v^2$, it is the message sent
  across the separator $\{\xi_m\}$ to a leaf bag $\{\xi_m\}$ carrying
  $\xi_m^2$. If $[0,2^m-1]$ is covered by $N_{\rm pc}$ intervals on each of
  which $V$ coincides with a quadratic polynomial, or if
  $V=\min_{l\le N_{\rm pc}}\varphi_l$ for quadratic polynomials $\varphi_l$,
  then $N_{\rm pc}\ge2^{m-1}$.
\end{enumerate}
By Theorem~\ref{thm:approx}, CT nevertheless solves this family to accuracy
$2^{-q}$ in time polynomial in $m+q$.
\end{proposition}

\begin{proof}
Part 1: $\xi_t$ lies in bags $t$ and $t+1$, and the endpoints $2^t-1$ have
$t$ bits.

Part 2. On the box $\xi_t\ge0$ and $z_t\ge0$, so
$\xi_{t-1}=(\xi_t-z_t-\rho_t)/2\le(\xi_t+\abs{\rho_t})/2$. By induction,
\[
0\le\xi_t\le2^{t-m}\xi_m+\sum_{j=t+1}^m2^{t-j}\abs{\rho_j} .
\]
The vector $(2^{t-m})_{t\le m}$ has squared norm at most $4/3$. The map
$y\mapsto(\sum_{j>t}2^{t-j}y_j)_t$ is the sum over $l\ge1$ of $2^{-l}$ times
the shift by $l$ places. Each shift has operator norm at most one, so the map
has operator norm at most $\sum_{l\ge1}2^{-l}=1$. Hence
$\norm\xi\le(2/\sqrt3)\xi_m+\norm\rho$ and
$\norm\xi^2\le\frac83\xi_m^2+2\norm\rho^2$. Also
$0\le z_t=\xi_t-2\xi_{t-1}-\rho_t\le\xi_t+\abs{\rho_t}$, so
$\norm z^2\le2\norm\xi^2+2\norm\rho^2$. Thus
$\norm\xi^2+\norm z^2\le3\norm\xi^2+2\norm\rho^2\le8\xi_m^2+8\norm\rho^2
\le8\Psi_m$.

Part 3. The second derivatives are $2+8=10$ for $\xi_t$, $t<m$; $2+2=4$ for
$\xi_m$; and $2-\frac14$ for $z_t$. Writing $\Psi_m$ as a sum of squares plus
$-\frac18\sum z_t^2$ plus linear terms gives
$\nabla^2\Psi_m\succeq-\frac14\sum_te_{z_t}e_{z_t}^{\T}\succeq-\frac14I$,
where $e_{z_t}$ is the unit vector of the coordinate $z_t$.

Part 4. Let $V_0(v)=V(v)-v^2
=\min\{\sum_t\rho_t^2+\frac18\sum_tz_t(1-z_t):\xi_m=v\}\ge0$. The minimum is
attained by compactness. Hence $V_0(v)=0$ exactly when some feasible point
has $\xi_m=v$, all $\rho_t=0$ and binary $z$. Then $\xi_t=2\xi_{t-1}+z_t$ and
$v=\sum_t2^{m-t}z_t\in\{0,\ldots,2^m-1\}$. Conversely, for an integer $j$ in
this range, its binary digits give $\xi_t=\lfloor j/2^{m-t}\rfloor\in
[0,2^t-1]$. So $V_0$ vanishes exactly at the $2^m$ integers of
$[0,2^m-1]$.

In the first representation, on an interval with nonempty interior,
$V_0$ coincides with a polynomial of degree at most two. It is not
identically zero, since $V_0>0$ off the integers. So it has at most two
zeros, and the interval contains at most two integers. A degenerate interval
contains at most one. Hence $2N_{\rm pc}\ge2^m$. In the second
representation, each integer $j$ has some $l$ with $\varphi_l(j)=V(j)$. Then
$\varphi_l-v^2\ge V_0\ge0$ on $[0,2^m-1]$ and vanishes at $j$. If
$\varphi_l-v^2\equiv0$, then $V_0\le0$, which is false; so $\varphi_l-v^2$
vanishes at most at two integers, and again $2N_{\rm pc}\ge2^m$.
\end{proof}

\subsection{Point growth cannot be replaced by set growth}
\label{sec:limits-setgrowth}

Sections~\ref{sec:constraints} and~\ref{sec:optsets} use set growth
(Definition~\ref{def:growth}). The next proposition shows that set growth
alone does not bound the grids that filtering produces.

\begin{proposition}[Set growth does not bound filtered grids]\label{lim:prop:setgrowth}
For $n\ge2$ let $F(x)=\sum_{i=1}^{n-1}(x_i-x_{i+1})^2$ on $X=[0,1]^n$, with
the path decomposition of bags $\{x_i,x_{i+1}\}$. The optimal set is the
diagonal $\mathcal S=\{t\mathbf1:t\in[0,1]\}$, $\OPT=0$, and
$F(x)\ge\frac2{n(n-1)}\dist(x,\mathcal S)^2$; for $n=2$,
$F=2\dist(x,\mathcal S)^2$. Every valid upper coordinate curvature satisfies
$L_i\ge2$, and with the smallest valid ones, $\kappa_S\le2n(n-1)$. Let the
corrections be $d_i(v)=L_iw_i(v)^2/8$ (Definition~\ref{def:corr}).
\begin{enumerate}[label=(\alph*)]
\item For every grid $G$ of $X$, every $j$ and every node $v\in G_j$, the
  min-marginal satisfies $m_j(v)\le-\frac{L_j}8w_j(v)^2$. In particular
  $\beta\le-\frac{L_j}8W_j^2$, where $W_j$ is the largest length of a grid
  interval of $G_j$.
\item The filtering step (Definition~\ref{def:filter}), applied to a grid
  of $X$ with a threshold $U\ge\OPT$, removes nothing. Hence every run of
  filtering with thresholds $U_j\ge\OPT$ that starts from a grid of $X$, in
  particular every run of TRIAL (Algorithm~\ref{alg:trial}) and hence of CT,
  uses grids of $X$ at all stages, and if it certifies accuracy
  $\varepsilon$, its last grid has at least $1+1/(2\sqrt\varepsilon)$ nodes
  in every coordinate.
\item In every valid path certificate (Definition~\ref{def:cert}) with
  $F(\hat x)-\beta\le\varepsilon$, every grid interval of $G^{(0)}_i$ that is
  not contained in $X^{(1)}_i$ (every grid interval of $G^{(0)}_i$ if $k=0$)
  has length at most $2\sqrt\varepsilon$.
\end{enumerate}
\end{proposition}

\begin{proof}
For $x\in X$, $\dist(x,\mathcal S)^2\le\sum_i(x_i-x_1)^2$, and
$(x_i-x_1)^2\le(i-1)F(x)$ by Cauchy--Schwarz, which gives the set-growth
bound. For $n=2$, $\dist(x,\mathcal S)^2=(x_1-x_2)^2/2$. The coordinate second
derivatives are $2$ for $x_1,x_n$ and $4$ otherwise, so every valid $L_i$ is
at least $2$; with the smallest valid ones, $L=\max_iL_i\le4$, and $L=2$ if
$n=2$, so $\kappa_S\le2n(n-1)$.

(a) Fix $j$ and $v\in G_j$, and put $v_j=v$. Choose the remaining coordinates
outward along the path. If $v_k$ has been chosen and $i=k\pm1$ is the next
index, let $[\alpha,\alpha']$ be a grid interval of $G_i$ containing $v_k$, of
length $\delta_i$, and let $v_i$ be its endpoint nearest to $v_k$. Such an
interval exists because $G_i$ contains $0$ and $1$. Then
$(v_k-v_i)^2\le\delta_i^2/4$ and $w_i(v_i)\ge\delta_i$. Every edge of the
path joins some $i\ne j$ to the index chosen just before it, so
\[
Q(v)=F(v)-D(v)\le\sum_{i\ne j}\frac{\delta_i^2}4-\frac{L_j}8w_j(v)^2
-\sum_{i\ne j}\frac{L_i}8\delta_i^2\le-\frac{L_j}8w_j(v)^2,
\]
because $L_i\ge2$. Hence $m_j(v)\le Q(v)\le-\frac{L_j}8w_j(v)^2$, and an
endpoint $v$ of a longest grid interval gives the bound on $\beta$.

(b) Let $G$ be a grid of $X$, let $[a,a']$ be a grid interval of $G_i$, and
let $t\in[a,a']$. The point $t\mathbf1$ is optimal and lies in $X$, so
Proposition~\ref{prop:cellwise}(c) at $x=t\mathbf1$ gives
$\min\{m_i(a),m_i(a')\}\le F(t\mathbf1)=0\le U$, and the interval is
retained. By induction, every stage of such a run has the box $X$. Suppose that the
last stage certifies accuracy $\varepsilon$, that is, $U-\beta\le\varepsilon$
with $\beta$ the corrected minimum of the last grid and $U\ge\OPT=0$. By (a)
and $L_j\ge2$, $W_j^2/4\le-\beta\le U-\beta\le\varepsilon$, so
$W_j\le2\sqrt\varepsilon$. Covering $[0,1]$ by intervals of length at most
$2\sqrt\varepsilon$ needs at least $1/(2\sqrt\varepsilon)$ of them, so
$\abs{G_j}\ge1+1/(2\sqrt\varepsilon)$.

(c) Since $F(\hat x)\ge\OPT=0$, $\beta\ge-\varepsilon$. The grid $G^{(0)}$ is
a grid of $X^{(0)}=X$. Let $[a,a']$ be a grid interval of $G^{(0)}_i$ of
length $\delta>0$ that is not contained in $X^{(1)}_i$. Both $a$ and $a'$ are
endpoints of it, so $w_i(a),w_i(a')\ge\delta$, and (a) gives
$\min\{m^{(0)}_i(a),m^{(0)}_i(a')\}\le-\delta^2/4$. Condition (C1) requires
this minimum to be at least $\beta\ge-\varepsilon$ (its second alternative
needs $w(J)=0$, that is, $\delta=0$), so $\delta\le2\sqrt\varepsilon$. If
$k=0$, (C2) and (a) give $-W_i^2/4\ge\min_{G^{(0)}}Q^{(0)}\ge\beta\ge-\varepsilon$
for every $i$.
\end{proof}

Part (b) covers the grids that CT returns. A general path certificate may
remove intervals whose endpoint min-marginals lie in $[\beta,0)$, and then
its last grid can be small; for such certificates, part (c) bounds only the
stage-$0$ grid. For
$n=2$ and $\varepsilon=1/50$, the uniform grid of step $1/5$ on $X$, followed by
the box $[0,1/5]^2$ with the grid $\{0,1/5\}^2$, $\beta=-1/50$ and $\hat x=0$,
is a valid path certificate of accuracy $\varepsilon$: all stage-$0$
min-marginals equal $-1/50$. Its last grid has two nodes per coordinate,
while (b) requires at least five for runs of filtering, and its removed
intervals have length $1/5\le2\sqrt\varepsilon$. We do not know whether every
valid path certificate for this family has $\varepsilon^{-\Omega(1)}$ nodes in
total. For finitely many optimal values per coordinate, unions of uniform
cells give grids with $O(r\sqrt{n\kappa_S})$ nodes per coordinate
(Theorem~\ref{thm:cells}). The proposition concerns a
continuum of optimal coordinate values. It is a limit of filtered grids, not
of the problem: here $F$ is a sum of squares.

\begin{proposition}[Value-oracle barrier for unknown set growth]\label{prop:oraclebarrier}
Fix $n\ge1$, $X=[0,1]^n$ and the coordinate-curvature bound $L=1$. Consider a
deterministic algorithm that queries exact values, gradients and Hessians of
an unknown $C^2$ function $f$ at adaptively chosen points and outputs either a
rational interval of width at most $\varepsilon$ containing $\min_Xf$, or a
point $\hat x$ with $f(\hat x)-\min_Xf\le\varepsilon$. Suppose it is correct
for every $C^2$ function $f$ with $\partial_{ii}f\le1$ on $X$ for all $i$ that
has set growth, $f(x)-\min_Xf\ge g_S\dist(x,\argmin_Xf)^2$ on $X$ for some
$g_S>0$. Then on $f\equiv0$, for which $\kappa_S=1$, it makes at least
$\bigl(1/(2\sqrt{6\varepsilon})-1\bigr)^n-1$ queries.
\end{proposition}

\begin{proof}
For $c\in X$ and $0<w\le1$ let
$f_{c,w}(x)=-\frac{w^2}6\bigl(1-\norm{x-c}^2/w^2\bigr)^3$ when
$\norm{x-c}\le w$ and $f_{c,w}(x)=0$ otherwise. In the variable $z=(x-c)/w$
with $\rho=\norm z^2$,
\[
\partial_if_{c,w}=w\,z_i(1-\rho)^2,\qquad
\partial_{ij}f_{c,w}=\delta_{ij}(1-\rho)^2-4z_iz_j(1-\rho).
\]
Value, gradient and Hessian vanish on $\rho=1$, so $f_{c,w}$ is $C^2$, and
$\partial_{ii}f_{c,w}=(1-\rho)(1-\rho-4z_i^2)\le1$. The minimizer $c$ is
unique with value $-w^2/6$. Inside the ball,
$f_{c,w}+\frac{w^2}6=\frac{w^2}6\bigl(1-(1-\rho)^3\bigr)\ge\frac{w^2}6\rho
=\frac{\norm{x-c}^2}6$, because $1-(1-\rho)^3-\rho=\rho(1-\rho)(2-\rho)\ge0$.
Outside it the gap is $w^2/6\ge\frac{w^2}{6n}\norm{x-c}^2$. Hence set growth
holds with $g_S=w^2/(6n)$. The zero function has $\argmin_Xf=X$ and satisfies
it with every $g_S$, in particular with $g_S=1$.

Run the algorithm on $f\equiv0$; let $Q$ be its number of queries and add
the output point, if any, to the query points. Suppose
$Q<\bigl(1/(2\sqrt{6\varepsilon})-1\bigr)^n-1$, and let
$M=\lfloor(Q+1)^{1/n}\rfloor+1$. Then $M^n>Q+1$ and $M<1/(2\sqrt{6\varepsilon})$.
Divide $X$ into $M^n$ closed subcubes of side $1/M$. The interior of one of
them contains none of the at most $Q+1$ points. Let $c$ be its center and
choose $w$ with $\sqrt{6\varepsilon}<w<1/(2M)$. The closed ball of radius $w$
lies in the open subcube, so every oracle answer for $f_{c,w}$ at the
recorded points equals the answer for $0$. By determinism the algorithm
produces the same transcript and output on $f_{c,w}$. An interval of width at
most $\varepsilon$ containing $0$ does not contain $-w^2/6<-\varepsilon$. An
output point outside the ball has gap $w^2/6>\varepsilon$ for $f_{c,w}$.
Either way the algorithm is wrong on an admissible function.
\end{proof}

CT uses only exact values and the curvature bounds, yet under point growth
its number of nodes per coordinate does not depend on the accuracy.
Proposition~\ref{prop:oraclebarrier} shows that no algorithm in that
information model, even with exact derivatives, can do the same under growth
towards an unknown set, even with $n=p=1$ and $\kappa_S=1$. The barrier
concerns information, not the input class: the functions $f_{c,w}$ are not
quadratics. It is the classical hidden-bump argument of information-based
complexity \cite{NemirovskiYudin1983}, adapted to set growth. An algorithm
for explicit quadratics must therefore exploit the coefficients, for example
through a certificate such as Lemma~\ref{lem:diagcert}, or the structure of
Section~\ref{sec:endpointset}.

\subsection{Product domains are essential}\label{sec:limits-constraints}

The next construction is the running-sum encoding of Subset Sum of Bienstock
and Mu\~noz \cite{BienstockMunoz2018}; see also
\cite[Example~1.1]{CifuentesParrilo2016}. What we add is a unique minimizer
and the growth constant, so that $\kappa=1$. An extra item $x_0$, whose size
is the target, keeps the feasible set nonempty.

\begin{proposition}[Local affine constraints]\label{lim:prop:constraints}
Let $a_0,a_1,\ldots,a_m$ be positive integers with $a_i\le a_0$. Consider
binary $x_0,\ldots,x_m$, continuous $y_0,\ldots,y_{m+1}\in[0,1]$, and the
constraints
\[
y_0=0,\qquad y_{m+1}=1,\qquad y_{i+1}=y_i+\frac{a_i}{a_0}x_i\quad(i=0,\ldots,m),
\]
with objective $F=2^{m+1}x_0+\sum_{i=1}^m2^{i-1}x_i$. Each constraint lies in
the bag $\{y_i,x_i,y_{i+1}\}$, and these bags form a path with $p=3$. The
feasible set is nonempty, the minimizer $z^*$ is unique, and
$F(z)-\OPT\ge\norm{z-z^*}^2/(2m+3)$ for every feasible $z$. Since $F$ is
affine, every $L>0$ is valid; with $L=1/(2m+3)$, $\kappa=1$. Moreover
$x_0^*=0$ if and only if some subset of $\{a_1,\ldots,a_m\}$ sums to $a_0$.
Consequently, unless $\mathrm P=\mathrm{NP}$, no algorithm solves such
problems in time $f(p,\kappa)\poly(I)$ for any function $f$, even when the
optimal value is required only within $1/2$.
\end{proposition}

\begin{proof}
Summing the recurrences gives $\sum_{i=0}^ma_ix_i=a_0$. The choice $x_0=1$,
$x_i=0$ ($i\ge1$) is feasible. Since all $a_i>0$, it is the only feasible
choice with $x_0=1$. A choice with $x_0=0$ is feasible exactly when its
items sum to $a_0$: its partial sums then lie in $[0,1]$. For each binary $x$
the constraints determine $y$. Distinct feasible points therefore have
distinct binary vectors and distinct integer objective values. The dummy
value $2^{m+1}$ exceeds every value $\le2^m-1$ of a genuine subset. Hence
the minimizer is unique, and $x_0^*=0$ exactly for yes-instances. All
$2m+3$ coordinates lie in $[0,1]$, so $\norm{z-z^*}^2\le2m+3$, while
$F(z)-\OPT\ge1$ for $z\ne z^*$. An approximation within $1/2$ separates
$\OPT\le2^m-1$ from $\OPT=2^{m+1}$. The data have polynomial binary length,
and Subset Sum is NP-complete \cite{GareyJohnson1979}.
\end{proof}

In the product setting an affine objective is solved exactly by endpoint
dynamic programming. The difficulty here comes entirely from the
constraints. Independent coordinate rounding, the basis of
Proposition~\ref{prop:cellwise}, does not preserve the running-sum equations.

After scaling $y$ by $a_0$, the instances of
Proposition~\ref{lim:prop:constraints} satisfy Definition~\ref{def:tu-model}
with a path constraint matrix, mesh unit $\eta=1$, box widths $s=a_0$ and two
values per discrete coordinate; the objective is affine, so every $\bar L>0$
is valid, and $\bar L\le g_S$ gives $\kappa_c=1$. By the discussion after
Theorem~\ref{thm:tu-approx}, the pseudopolynomial level-$0$ grid of
TU-GRID (Algorithm~\ref{alg:tugrid}) therefore cannot be removed unless
P${}={}$NP.

\subsection{Local corrections and local moments}\label{sec:limits-local}

Proposition~\ref{prop:star} shows that a correction charged only to the
coordinates of one bag, while the other coordinates are optimized on the
grid, cannot be valid for any fixed allowance: the missing error grows with
the number of outside coordinates and depends on the separator value. The
valid local alternative is exact or certified conditional recourse
(Theorem~\ref{thm:cr-filter}).

A different route to a local error bound keeps the exact positive curvature
and discretizes only the negative part: for a continuous box quadratic with
$\nabla^2F\succeq-\nu I$ one would like an error bound in which $L$ is
replaced by $\nu$. Such a bound holds for relaxations that condition all
coordinates on one global cell (Lemma~\ref{lim:lem:mixture}), but a global
cell index can take $K^n$ values. Sparse relaxations keep exact information
only on bags and match moments across separators, cell by cell
(Definition~\ref{lim:def:moments}). Proposition~\ref{lim:prop:moments} shows
that this does not suffice for any fixed order $k$ of the matched moments: for
every $k$ there are instances with bag size three, a unique minimizer,
$\nabla^2F\succeq-2I$ and a growth constant that depends only on $k$, whose
order-$k$ relaxations have feasible points more than $1/8$ below $\OPT$ that
occupy only cells of arbitrarily small length. A $\nu$-based local error
bound must therefore use information beyond finite-order separator moments,
for example inequalities that couple both sides of a separator.
Appendix~\ref{app:moments} gives the definitions, the construction and the
proofs.
